%! TEX root = ../main.tex
\documentclass[main]{subfiles}

% 全体的にCMT_scaleの記述を修正
\begin{document}
\section{はじめに}
日本における少子高齢化の進行に伴い，医療・介護資源の不足は喫緊の課題である．
この解決策の一つとして，情報技術を活用して高齢者や障がい者の生活支援を行う Ambient Assisted Living (AAL) が挙げられる．
AALは在宅での自立生活を促進することで，医療や介護で発生する労働時間や社会保障費などの社会的コストを抑制すると同時に，
被介護者の生活の質（Quality of Life; QoL）向上に寄与することを目指すものである\cite{blackman2016ambient}．\par

AALの要素技術である行動認識は，センサを用いて居住者の行動を認識し，異常行動の検知や健康状態の把握を行う．
従来の研究では，主にウェアラブルセンサ，カメラ，あるいはその他の環境設置型センサが用いられてきた．\par

ウェアラブルセンサは，対象者の状態を常時取得でき，転倒などの異常を即時に検知できる．
しかし，高齢者の認知機能低下に伴う装着忘れや充電等の管理不備のリスクが伴う．
加えて，常時機器を身につけることによる身体的な不快感や拘束感のみならず，被監視による心理的圧迫感が課題となる\cite{guerra2023ambient}． \par

一方，カメラは情報量が豊富で具体的な行動認識が容易な上，身体への装着も不要である．
しかし，マイクロフォンと比較して導入コストが高くなりがちであり，プライバシーの観点から被介護者の心理的抵抗感が強い場合もある．
また，撮影範囲の死角や，夜間の暗所において認識精度が低下する点も課題である\cite{guerra2023ambient,bouchabou2021survey}．\par

その他の非接触な手法として，家具や床に設置する圧力センサや，無線通信波の変動を利用して行動を検知する環境センサなどを利用する方法がある\cite{guerra2023ambient}．
これらは身体的負担がない点で優れるが，本研究ではこれらに加え，聴覚情報が持つ「環境音」や「発話」などの多面的な情報に着目する．\par

マイクロフォンは設置が容易でありながら，死角の影響を受けにくく，かつカメラよりもプライバシー受容性が高い．
特に，カメラの設置が倫理的に難しい浴室や洗面所においては，その有用性は高いと考えられる\cite{hyun2023sound}．
例えば，ヒートショックによる転倒や溺水が発生した際，流水音の継続や転倒時の衝撃音といった特定の音響イベントを捉えることで，
視覚的なプライバシーを侵害することなく異常を検知できる可能性がある．\par

しかし，実環境におけるマイクロフォンの運用には課題も残る．
換気扇の稼働音や生活雑音などの環境ノイズ，あるいは部屋の形状に起因する残響特性の違いが，
認識精度を低下させる要因となるためである\cite{mesaros2021sound}． 

本研究の長期的展望としては，マイクロフォンを含む複数のセンサを統合し，多様な文脈を考慮した柔軟な行動認識システムの構築を目指している．
本論文ではその初期検討として，実環境特有のノイズや環境要因に対し，マイクロフォンによる行動認識精度を向上させる手法について検討する．\par

他方， 音響信号から事象の種類と発生時刻を特定する音響イベント検出（Sound Event Detection; SED）の研究が活発化している．
SEDは， 音響イベントのクラス識別および時間区間検出を同時に行う技術であり， 本論文では音による行動認識の手段としてこれを位置付ける．
本タスクでは，10秒間の音声データを「クリップ」，それをさらに短時間で分割した単位を「フレーム」と定義し，それらを入力として受け取り，イベント種別及び各イベントの開始・終了時刻を予測する．\par

家庭内環境向けSEDのベンチマークとして， Detection and Classification of Acoustic Scenes and Events (DCASE) チャレンジが知られている． 
本研究では， DCASE 2024 Task 4 (DCASE 2024) の評価基準モデルを参考実装（ベースライン）\footnote{\url{https://github.com/DCASE-REPO/DESED_task/tree/master/recipes/dcase2024_task4_baseline}}とし， 
評価用データセットとしてDomestic Environment Sound Event Detection Dataset (DESED)，
および MAESTRO Real を使用する\cite{turpault2019sound,serizel2020sound,martin2023strong}．\par

これらデータセットのラベルには種類がある．弱ラベルはイベント種別のみを，強ラベルは開始・終了時刻も含み，それぞれクリップとフレームに対応する．
本タスクではクリップでイベントの有無を，フレームではイベント種別に加えて，その開始・終了時刻を学習・予測する． \par

データセットについて，DESEDは家庭内の音響イベント検出に向けて，実環境で録音した家庭内の生活音と，そのような環境を再現するために合成された音が含まれる．
音声データは10秒間のモノラルオーディオで，サンプリングレート 16kHz または 44.1kHzとなっている．
% モノラル > チャネル数が1
MAESTRO Realは実環境の$3-5$分の音が含まれる．ラベルは0か1ではなく，0.8などの連続値 (ソフトラベル) で，強ラベルが付与されている．\par


DCASE 2024では， 大規模な未ラベルデータを用いた複数の事前学習済みモデルの活用が精度向上に寄与し， 上位成績を収めた\cite{schmid2024improving}．
 しかし， これらの手法はモデルの大規模化を招き， パラメータ数が著しく増大するという課題がある． 
 実環境への実装を想定した場合， モデルのパラメータ増加は望ましくない． 
 そこで本研究では， モデルサイズを維持しつつ精度を向上させることを目的に，検討を行った．なお，将来的にはモデルサイズの縮小を検討する． 
 はじめに半教師あり学習の改良を提案し，さらにデータ拡張によるドメイン汎化と後処理の適正化の3つを統合するアプローチを提案する．\par

DCASE 2024のベースラインモデルでは，ラベル欠損に対応する半教師あり学習手法としてMean Teacherが採用されている．
Mean Teacherは，通常のモデルを「生徒」，生徒モデルの重みパラメータの指数移動平均（Exponential Moving Average; EMA）を適用したものを「教師」と定義する．
教師モデルがラベルなしデータに対して行う予測を擬似的な目標値とし，生徒モデルの予測との一貫性を高めることで学習を行う\cite{tarvainen2017mean}．\par

 本研究では，MeanTeacherの疑似ラベルに閾値を設け，信頼できるラベルのみ学習に用いることで精度向上を図るConfident Mean Teacher (CMT) を採用する\cite{xiao2023multi}．
 CMTは確率値が0.5より大きい予測を信頼できると考えるが，閾値未満の確率値を持つ予測は損失に寄与せず，イベントがないことを学習するのが難しいと思われる．
 追試の結果，イベントの開始や終了時刻の予測精度（時間的な精度）が低下した一方で，イベントの分類予測では精度の向上が見られた．
 なお，CMT及びその検証を行っている論文では精度の向上が報告されているが，いずれも実装が非公開で，検証が難しい\cite{11036061}．\par 

 当環境で見られた精度低下について，フレーム単位でのイベントがない情報を学習できていないから，という仮説を立て，検証した．
 具体的には，確率値0.7以上の予測を正例，0.3以下の予測を負例とし，正例と負例の両方で学習を行った．確率値が0または1に近づくほどモデルが自信を持ち，信頼できると考える．
 CMTに比べ，あまり自信のない確率値0.5付近の予測を使わず，より安定した精度が得られると共に，負例の利用による時間的精度の向上が期待される．
 また，正例と負例で損失の量が異なる場合を考え，より少ない方に調整することも検討した．

次に，ドメイン汎化について述べる．ここでドメインとは，同一のクラスラベルであっても，データが観測された環境や条件によって生じる特徴分布の違いを指す．
例えば，画像認識において実写とイラストが異なるドメインとして扱われるのと同様に，音声においても録音環境や背景雑音の違いにより，同じ音響イベントであっても周波数特性などの特徴が大きく異なる場合がある．
モデルがこうした環境依存の情報に過剰適合することを防ぎ，未知の環境でもロバストな推論を可能にすることがドメイン汎化の目的である．\par

MixStyle\cite{zhou2021domain}はドメイン汎化に有効なデータ拡張手法の一つである．
先行研究\cite{xiao2024fmsg}において周波数領域の注意機構との統合が示唆されているが，その具体的なアーキテクチャについては詳述されていない．
そこで本研究では，軽量な畳み込みニューラルネットワーク (Convolutional Neural Network; CNN)ベースの手法を検討した． 
% これが本当に僕の工夫で良いのか? 確認が必要

提案手法では，まずCNNを用いて周波数方向の重要度に基づき特徴量を重み付けする．
次に，その特徴量に対してMixStyleを適用して統計量を混合し，最後に入力特徴量と加算（残差接続）した． 
これにより，タスクに重要な周波数情報を保持したままスタイルのみを変更することが可能となり，かつ残差接続によって特徴の消失を防ぎながら，効果的なドメイン汎化を実現する．
なお，\cite{kim2022domain}より，音声におけるドメイン情報は主に周波数領域の統計量に関連することが示唆されている．\par

最後に，SEDにおける後処理の改良手法としてSound Event Bounding Boxes（SEBBs）を導入する\cite{ebbers2024sound}．
SEDでは，一定時間を切り出した「クリップ」と，それをさらに細分化した「フレーム」の双方において，正しい予測が行われているかを評価する必要がある．
従来のDCASE 2024のベースライン等では，クラスごとにパラメータを調整したMedian Filter（中央値フィルタ）が一般的に用いられてきた．
本研究ではより柔軟なパラメータ調節が可能であるSEBBsを用いることで，評価スコアの改善を目指す．\par

以上の手法について，評価実験を行った．
評価指標として，DCASE 2024に基づき，DESEDの時間的な精度とイベントの分類精度，MAESTRO Realの総合的な予測の精度を見る．
本研究では，将来的に音響イベント検出における時間精度の向上とイベント分類の最適化を分離し，段階的に扱うことを想定している．
現段階ではクラス分類性能を重視するが，モデルの総合的な性能を評価するため，時間的な精度も併せて報告する．\par

実験の結果，CMTとMixStyle，SEBBsの三手法の組み合わせは，ベースラインより精度が向上した．
また，CMTの改良手法はCMTの精度を上回った．一方で，MixStyleなどと併用した際は，CMTと改良手法の間でほとんど差は見られなかった．\par

本論文は次のように構成される．2節では本研究が参考にしている実装を簡単に紹介し，提案手法について述べる．3節では実験で用いられる指標と，実験内容を述べる．
4節は実験の結果を整理し，考察を行う．5節は本論文の成果をまとめ，今後の課題を述べる．
\end{document}